% Response letter written from response-JL.docx (co-authors' version, 2026-11-19).
% Content follows the docx; layout follows the texpage-sync response-letter style.
% Compile with pdfLaTeX, twice. Figures come from ../letter_figs via the figures/ folder.
\documentclass[11pt,a4paper]{article}
\usepackage[utf8]{inputenc}
\usepackage[T1]{fontenc}
\usepackage{lmodern}
\usepackage[left=2.3cm,right=2.3cm,top=2.3cm,bottom=2.3cm]{geometry}
\usepackage{xcolor}
\usepackage{textcomp}
\usepackage{amsmath}
\usepackage{graphicx}
\usepackage[hidelinks]{hyperref}
\usepackage{parskip}
\usepackage{caption}
\usepackage{booktabs}
\usepackage{etoolbox}
\captionsetup{font=small,labelfont=bf,labelsep=period}
\renewcommand{\thefigure}{R\arabic{figure}}
\renewcommand{\thetable}{R\arabic{table}}
% \rfig{width}{file in figures/}{caption}: not a float, it stays where it is called. The minipage keeps image and caption on one page.
\newcommand{\rfig}[3]{\par\medskip\noindent\begin{minipage}{\linewidth}\centering\includegraphics[width=#1\linewidth,height=0.6\textheight,keepaspectratio]{figures/#2}\captionof{figure}{#3}\end{minipage}\par\medskip}

\definecolor{chgcol}{RGB}{0,70,140}
\pretocmd{\section}{\color{black}}{}{}
% Greek letters and math signs pasted from reviewer comments.
\DeclareUnicodeCharacter{03B4}{\ensuremath{\delta}}
\DeclareUnicodeCharacter{0394}{\ensuremath{\Delta}}
\DeclareUnicodeCharacter{03B1}{\ensuremath{\alpha}}
\DeclareUnicodeCharacter{03B2}{\ensuremath{\beta}}
\DeclareUnicodeCharacter{03C3}{\ensuremath{\sigma}}
\DeclareUnicodeCharacter{03BC}{\ensuremath{\mu}}
\DeclareUnicodeCharacter{2248}{\ensuremath{\approx}}
\DeclareUnicodeCharacter{2264}{\ensuremath{\leq}}
\DeclareUnicodeCharacter{2265}{\ensuremath{\geq}}
\DeclareUnicodeCharacter{2212}{\ensuremath{-}}
\newcommand{\reviewer}[1]{\section*{#1}}
% Comments are pasted verbatim in black: % # & _ ~ inside the second argument are printed as typed.
% The first argument is a tag for the source only and is not printed.
\newcommand{\comment}[1]{\par\medskip\noindent\color{black}\begingroup\catcode`\%=12 \catcode`\#=12 \catcode`\&=12 \catcode`\_=12 \catcode`\~=12 \commentbody}
\newcommand{\commentbody}[1]{#1\endgroup\par\smallskip}
% Responses are blue until the next comment or heading; quotes from the manuscript are blue italics.
\newcommand{\response}{\par\noindent\color{chgcol}\textbf{Response.}\ }
\newenvironment{change}{\par\smallskip\color{chgcol}\small\itshape\leftskip=1.2em\rightskip=1.2em\noindent\ignorespaces}{\par\smallskip}
\emergencystretch=2em

% ---- additions for this letter (aabw_5exps, Nature Communications revision) ----
\DeclareUnicodeCharacter{03B3}{\ensuremath{\gamma}}
% Yellow highlight kept from the co-authors' Word file. ulem (in every TeX Live) instead of soul,
% because soul before v3.0 breaks on UTF-8 characters under pdfLaTeX.
\usepackage[normalem]{ulem}
\DeclareRobustCommand{\hl}[1]{\bgroup\markoverwith{\textcolor{yellow}{\rule[-.5ex]{2pt}{2.6ex}}}\ULon{#1}}
% Word comment from the co-authors' file: number, author and date, the text it was attached to, the comment.
\DeclareRobustCommand{\docxcomment}[4]{{\upshape\sffamily\footnotesize\color{red}[Comment #1 (#2) on ``#3'': #4]}}

\begin{document}

\begin{center}
{\Large\bfseries Summary of Responses to Reviewers}\\[4pt]
{\large Glacial--interglacial regime shift in Southern Ocean dense water formation}\\[2pt]
Xiaoxu Shi, Jiping Liu, Chaoyuan Yang, Gerrit Lohmann
\end{center}

\bigskip
%=====================================================================

% docx heading: Response to comments by Reviewer 1

\reviewer{Reviewer 1}

\response We thank the reviewer for the constructive general comments and a thorough list of specific comments.

\comment{R1}{The manuscript presents a comprehensive investigation of Southern Ocean dense water formation across a range of glacial and interglacial climates using the AWI-ESM2 Earth system model. By applying a water mass transformation (WMT) framework, the authors successfully separate the relative roles of thermal and haline buoyancy forcing and demonstrate a compelling regime shift between interglacial and glacial climates. The manuscript is well structured, the analyses are technically rigorous, and the decomposition of the different buoyancy contributions provides valuable physical insight. The use of WMT diagnostics allows the underlying mechanisms to be quantified rather than inferred from changes in circulation alone. Overall, I found the study scientifically interesting and suitable for publication after major revisions. My comments below are primarily intended to strengthen the robustness of the conclusions and improve the clarity of the presentation.}

\comment{R1}{General comments}

\comment{R1}{Model evaluation and confidence in simulated dense water formation}

\comment{R1}{The conclusions of this study rely entirely on a single Earth system model. While I do not consider this to be a major limitation by itself, I think the manuscript would benefit from a more thorough discussion of how well AWI-ESM2 represents present-day Southern Ocean dense water formation and Antarctic Bottom Water formation. The manuscript briefly states that the PI simulation produces comparable WMT magnitudes to observational estimates, but Figure A4 only shows the model results, making it difficult for the reader to evaluate this statement. Given the importance of the model for all subsequent paleoclimate interpretations, I encourage the authors to provide stronger evidence supporting its performance. This could include a direct comparison with observational WMT estimates, or a more detailed discussion of known model strengths and biases in Southern Ocean stratification, sea ice, mixed-layer depth, and Antarctic Bottom Water formation. Even if such comparisons already exist in previous AWI-ESM2 publications, briefly summarising them here would substantially increase confidence in the conclusions.}

\response We agree with the reviewer. This concern also overlaps with those from Reviewers 2 and 3.

We now place the water mass transformation (WMT) of our pre-industrial simulation beside the observation-based estimate of Pellichero et al. (2018). The estimate of Pellichero et al. (2018) is based on Argo floats, ship-based observations, and instrumented marine mammals (the only platform that samples under winter sea ice), the ice edge derived from SMMR and SSM/I passive microwave data for 2000-2015, and the geostrophic term derived from AVISO mean dynamic topography for 1993-2012. We cannot reproduce that calculation, because it rests on a seal-tag mixed layer database those the authors assembled and we do not have. We have therefore compared our modelled WMT against their published figure.

To match the published figure of Pellichero et al. (2018), we recomputed our WMT over the region enclosed by the 15\% September sea ice concentration, and took the annual mean (used in Pellichero et al., 2018) instead of the austral winter mean used in our manuscript.

\rfig{0.9}{../letter_figs/figR20_wmt_merged.png}{New analysis. Transformation rate against neutral density γ, in three comparisons. Top row: annual mean for the seasonal sea ice zone, with all terms present: (a) our pre-industrial simulation state, the total is decomposed into the heat, freshwater and sea ice contributions, and (b) the observational estimate of Pellichero et al. (2018), their Fig. 3a, \hl{reproduced} for comparison. Pink shading in (a) marks where the heat term carries the total and blue shading where the freshwater term does, classified point by point along the drawn curve rather than by density bin, which places the thermally driven band at γ = 27.10–27.68 kg m\textsuperscript{-3}. Where the same band is quoted from the binned diagnostic of \hl{Fig. R2.3} it appears as γ = 27.33–27.59 kg m\textsuperscript{-3}, the single bin that falls inside it. Middle row: the same domain and season against the ORAS5 reanalysis for 1999-2018: (c) our pre-industrial simulation state with the sea ice contribution removed and (d) ORAS5, which carries heat and evaporation minus precipitation minus runoff only. Bottom row: the domain and season our paper uses, winter mean (JJA) for south of 60°S. (e) ours with the sea ice contribution removed and (f) ORAS5. Rows two and three are unshaded, because with the sea ice term absent from both sides, a heat-against-freshwater classification would describe an incomplete freshwater flux and read as a physical result rather than a missing field. Our simulated surface water is lighter than observed, so equal γ does not identify the same water mass across panels.\label{fR:wmt}}

As seen in Fig.~\ref{fR:wmt}a-b, the two figures share their structures. Both show buoyancy gain at the lighter classes, a zero crossing, and densification at the heavier classes, with the freshwater term carrying the light end and both curves converging towards zero at the dense end. Freshwater component dominate the light water class (< 27.1 kg/m3) and the dense end (> 27.67 kg/m3), shown by the blue shading in Fig.~\ref{fR:wmt}a, while the thermal signature is confined to a band at intermediate water densities, the pink shading in Fig.~\ref{fR:wmt}a. In our pre-industrial state, WMT reaches a minimum of -9.2 Sv, crosses zero, and peaks at +6.6 Sv; Pellichero et al. (2018) reported a minimum near -22 Sv and a densification peak of 5±5 Sv. The sign structure and the order of magnitude of the densification branch therefore agree.

Two differences should be noted. Our light-end buoyancy gain is weaker than theirs, and our curve sits more at lighter densities, because the simulated Southern Ocean surface is fresher than observed; both follow from the same bias, which we quantify under model limitations below.

It should be noted that the difference could be definitional, and it sets a floor on how closely the two can ever agree. Pellichero et al. (2018) solved the mixed layer heat and salt budgets for the surface fluxes as the unknown, using the observed seasonal cycle of mixed layer temperature, salinity and depth together with Ekman and geostrophic advection and entrainment. Thus, what Pellichero et al. (2018) solved for is the surface flux together with the diapycnal mixing term, because the mixed layer budget cannot separate them, whereas our model diagnoses the surface flux alone. \docxcomment{1}{Microsoft Office User, 2026-11-05}{alone.}{Also mention that although Pellichero et al. (2018) combined Argo floats, ship-based observations, and instrumented marine mammals to obtained WMT in the ice coved Southern Ocean, the data still have limited coverage and uncertainties.} Some discrepancy is therefore expected before any model bias is considered, and the agreement in shape should be read in that light.

The reviewer, and Reviewer 3, also asked whether we could compute an observational transformation from a reanalysis rather than citing a published one. A transformation budget needs the net surface heat flux, the net surface freshwater flux, and the sea ice thermodynamic salt flux. We checked the Copernicus ocean reanalyses, namely GLORYS12, the GREP ensemble and the analysis-forecast product, which carry temperature, salinity, velocity and sea ice state but no surface flux fields at all. ORAS5 does carry a net heat flux and a net water flux, but excluding the sea ice term. Therefore, we calculate transformation based on ORAS5, and for consistency, we also remove the sea ice term from our model, then compare our model results with ORAS5.

The results are shown in Fig.~\ref{fR:wmt}c-f. For annual mean over the sea ice zone, panels (c) and (d), the agreement is close. Our densification branch peaks at +5.0 Sv and the ORAS5 branch at +4.4 Sv, one density class apart at γ 27.49 against 27.28, and both show the light-end buoyancy gain, ours reaching a minimum of -8.0 Sv against -12.9 Sv in the reanalysis. The sign structure, the position of the zero crossing and the magnitude of the densification peak therefore agree on a domain neither calculation was tuned to.

For the domain and season our paper uses, south of 60°S in JJA, panels (e) and (f), the two agree in shape but not in amplitude. Both are positive almost everywhere, since in winter the whole sector is losing buoyancy, and both peak in the same density class at γ 27.28 and 27.08 respectively. Ours is the stronger, +62.9 Sv against +36.8 Sv. We do not attribute that gap to the reanalysis being correct. Our winter mixed layers south of 60°S are deeper and more extensive than in ORAS5, which exposes more water to the atmosphere and so supports more transformation, and the same bias is visible directly in the mixed layer depth comparison shown later in this response. We therefore read the winter comparison as agreement in sign, shape and density class, with an amplitude difference that follows from a mixed layer bias we document rather than from the diagnostic itself.

In the densest classes, above γ 27.7, neither curve sustains much transformation once the sea ice term is removed from both sides, which is the expected outcome. That is the class where brine rejection does the work, and it is exactly the term ORAS5 cannot supply. With the term restored, our own dense-class transformation is positive, which is the result the paper rests on and the one no available reanalysis can be used to check.

Two caveats bound what this comparison can support. ORAS5 is a model constrained by observations rather than a measurement, \hl{and the two fields the calculation needs are the two that are least constrained}, since no observing system measures an air-sea flux directly; its heat flux is computed from ERA-Interim forcing and then adjusted by sea surface temperature restoring, not fitted to any flux observation. Pellichero et al. (2018) rejected this approach for this region, stating that ocean buoyancy fluxes are commonly obtained from atmospheric reanalysis datasets but that ``in the Southern Ocean such reanalysis products are poorly constrained, and essentially unusable in the seasonally sea-ice covered region.'' We therefore present the ORAS5 comparison as a consistency check on the diagnostic.

Another point from the reviewer is about the discussion of known model strengths and biases in Southern Ocean stratification, sea ice, mixed-layer depth, and Antarctic Bottom Water formation. Since the second last comment of the reviewer specifically deals with ``Model limitations'', we refer to our response to that comment.

\comment{R1}{Boundary conditions}

\comment{R1}{The description of the prescribed boundary conditions for the paleoclimate simulations could be expanded. In particular, the glacial simulations use the GLAC-1D ice-sheet reconstruction, but it is not entirely clear how the ice-sheet changes are implemented in the model (e.g., topography, land-sea mask, ice-sheet extent, bathymetry, or other boundary conditions). As Antarctic dense water formation is highly sensitive to such factors, providing additional details on the experimental setup would help readers assess the robustness of the simulations. This is not a request for additional simulations, but rather for a clearer discussion of the robustness of the conclusions and the extent to which they may depend on the chosen model configuration and forcing datasets. A summary table listing the prescribed greenhouse gas concentrations, orbital parameters, and ice-sheet boundary conditions for each experiment would also improve the readability of the Methods section.}

\response Following the reviewer’s suggestion, we provided a summary table (Table R1, manuscript Table 1), which lists greenhouse gas concentrations, orbital parameters, and ice sheet configuration for all five AWI-ESM2 model experiments, and the values are those actually used in the model runs rather than the protocol defaults.

\begin{table}[h!]
\caption{Boundary conditions prescribed in the five equilibrium simulations. Orbital parameters for PI are those of 1850 CE. Greenhouse gas concentrations are volume mixing ratios.}
\label{tR:bc}
\centering\footnotesize\setlength{\tabcolsep}{3.5pt}
\begin{tabular}{lccccccc}
\toprule
 & CO\textsubscript{2} (ppm) & CH\textsubscript{4} (ppb) & N\textsubscript{2}O (ppb) & Eccentricity & Obliquity (°) & Perihelion (°) & Ice sheet / bathymetry \\
\midrule
PI & 284.3 & 808.2 & 273.0 & 1850 CE & 1850 CE & 1850 CE & modern \\
MH & 264.4 & 597.0 & 262.0 & 0.018682 & 24.105 & 180.87 & modern \\
LIG & 275.0 & 685.0 & 255.0 & 0.039378 & 24.040 & 95.41 & modern \\
LGM & 190.0 & 375.0 & 200.0 & 0.018994 & 22.949 & 294.42 & GLAC-1D 21 ka \\
MIS3 & 210.5 & 556.2 & 247.4 & 0.013676 & 23.2591 & 25.99 & GLAC-1D 38 ka \\
\bottomrule
\end{tabular}
\end{table}

We also expanded the text for describing the boundary conditions used in the experimental design, especially the implementation of the GLAC-1D ice-sheet reconstruction. Now reads:

\begin{change}
``We perform five equilibrium simulations spanning distinct climate states including three interglacial periods, i.e., pre-industrial (PI), mid-Holocene (MH), and last interglacial (LIG), and two glacial periods, i.e., the Last Glacial Maximum (LGM) and Marine Isotope Stage 3 (MIS3). The boundary conditions are configured following the criteria of PMIP4 and are summarised in Table 1. Orbital parameters are calculated following \textit{(Berger, 1977)}, and the greenhouse gas concentrations are taken from multi-archive reconstructions from ice core records and recent measurements of firn air and atmospheric samples. For LGM and MIS3, we fix the boundary conditions at 21 ka and 38 ka respectively, with the topography and ice-sheet properties deriving from the GLAC1D reconstruction. The GLAC-1D reconstruction enters the model through several boundary fields at once, and it is clearest to list them by the model component that uses them. Specifically, the atmospheric model component receives the ice sheet extent and surface elevation as orography, together with the land-sea mask implied by the lowered sea level. The land surface model component receives the same topography, from which river routing and vegetation distribution are derived. The ocean model component receives bathymetry and a modified coastline. Because the ocean component uses an unstructured mesh, a changed coastline cannot be imposed on the existing grid and a dedicated mesh must be generated for each glacial period. We use meshes constructed for 21 ka and 38 ka respectively, while PI, MH and LIG share the modern-geometry mesh. Global mean ocean salinity is raised in the glacial experiments to account for the water stored in the ice sheets. No prescribed ice sheet meltwater flux is applied in any of the equilibrium simulations, and the model does not include ice shelf cavities, so basal melt beneath floating ice shelves is not represented.''
\end{change}

\comment{R1}{Clarifying several interpretations}

\comment{R1}{Most interpretations are convincing, although a few statements could be presented more carefully or supported more explicitly. For example, some descriptions of the role of the Southern Annular Mode appear stronger than the presented evidence suggests, and several discussions would benefit from more direct references to the relevant figure panels.}

\response We have moderated the Southern Annular Mode claims. The most overstated passage was the suggestion that the simulated coupling implies a route to marine proxy reconstruction; we have replaced it with a more careful statement noting that the accumulation rates and chronological precision required are rarely available around the Antarctic margin, a point Reviewer 3 also raised. We have also removed the claim that the circumpolar picture holds uniformly and now state explicitly where it does not, in particular the Ross Sea. Panel-level figure references have been added throughout, as requested in this comment and in comment 21.

In the revised manuscript the moderated passage now reads:

\begin{change}
``Whether this persistent coupling could be exploited for proxy reconstruction is less clear than we previously suggested. Recovering an interannual-to-decadal signal from marine sediments would require accumulation rates and chronological precision that are rarely available around the Antarctic margin, where sedimentation is slow and bioturbation mixes the record over intervals far longer than the variability of interest. We therefore present the persistence of the coupling as a property of the simulated system rather than as a practical route to a marine Southern Annular Mode reconstruction.''
\end{change}

\comment{R1}{Model limitations}

\comment{R1}{Model limitations: The manuscript convincingly demonstrates that sea ice drives the transition from a thermally dominated to a haline-dominated dense water formation regime. However, the discussion would benefit from a brief consideration of the known biases and limitations of AWI-ESM2, particularly regarding Southern Ocean sea ice, mixed-layer depths, and Antarctic Bottom Water formation. Discussing how these model characteristics may influence the proposed mechanism would help readers assess the robustness and broader applicability of the conclusions.}

\response This comment asks specifically about model bias, where the general comment above asked about evaluation, so we treat the two comments separately. We take the three quantities the reviewer names in turn, and in each case compare our pre-industrial simulation against observations.

Thanks for the comment. We now state the model limitation explicitly. We diagnosed winter mixed layer depth (Figure~\ref{fR:mld}), reflecting where the model actually forms its dense water. The broad pattern shown in the ORAS5 reanalysis is reproduced by our pre-industrial simulation. Both the model and the reanalysis keep shallow mixed layers over most of the open Southern Ocean, deepen towards the Antarctic margin, and place the deepest water of the sector in the Atlantic sector. The obvious difference is in the Weddell gyre interior. In our model simulation, the Weddell sector accounts for 70\% of the area with a winter mixed layer deeper than 400 m south of 55°S, centred near 68°S in the interior of the gyre (see Figure~\ref{fR:mld}a). This is open-ocean convection, not the shelf overflow pathway that operates in the real ocean. We note that this is a limitation shared across the current model generation. Heuzé (2021) found that almost all CMIP6 models form deep and bottom water by open-ocean convection that is too deep, too frequent, or too extensive, with no model showing convincing shelf export in the Weddell Sea. We also now state that this configuration has no ice shelf cavities, so basal melt beneath floating ice is absent by construction.

\rfig{0.9}{../letter_figs/figR17_mld_pi_vs_obs.png}{New analysis. Winter mixed layer depth south of 60°S. (a) The pre-industrial simulation. (b) The ORAS5 ocean reanalysis, JJA climatology over 1990-2009. The large-scale pattern agrees: both are shallow over most of the open Southern Ocean, both deepen towards the Antarctic margin, and both show the deepest water of the sector in the Atlantic sector. The difference is where it matters for this study. The model places a large patch deeper than 400 m in the interior of the Weddell gyre, covering 1.9 per cent of the area south of 60°S against 0.0 per cent in the reanalysis, and its Weddell sector mean is 128 m against 71 m. Area-weighted over the whole sector the model gives 100 m against 78 m, so the bias is concentrated in the convecting patch rather than spread evenly. We note that the two use different mixed layer definitions, a 0.01 kg m\textsuperscript{-3} density threshold in ORAS5 and a bulk Richardson number criterion in the model, so the point-by-point difference is not purely model error.\label{fR:mld}}

\hl{Add a winter sea ice cover plot, model vs. satellite, and discuss the similarity and difference.}

In addition, the published evaluation reports \docxcomment{2}{Microsoft Office User, 2026-11-05}{the published evaluation reports}{Provide reference!} that Antarctic sea ice in this configuration is thiner \hl{compared to ?}, especially in the central Weddell Sea in September, and that the Southern Ocean has a warm bias of 1.43 K. The abyssal overturning cell is reported at about 10 Sv, comparable to our pre-industrial value. That cell is produced by the wrong pathway in the Weddell Sea, as described above, so we present our results as a statement about how the partitioning of surface buoyancy forcing responds to changing ice cover, and explicitly not as a quantitative reconstruction of bottom water formation rates.

Taken together these are systematic biases. That matters for how the results should be read. Because the same bias is present in all five experiments, it largely cancels in the differences between them, which is what the paper reports. We make no claim about the absolute transformation rate or the absolute density at which it occurs; we claim that the balance between the thermal and haline pathways shifts between climate states, and a bias common to every state cannot generate that shift.

The corresponding passage added to the Discussion reads:

\begin{change}
``Independent evaluations of this model provide context for these limitations, and we summarise the relevant points here rather than leaving them to be sought in the model description literature.\docxcomment{3}{Microsoft Office User, 2026-11-05}{literature.}{Also add a few more sentences from the above evaluation of mixed layer depth and sea ice cover here.} Sidorenko et al. (2019) evaluate the coupled configuration against observations and report that the abyssal overturning cell associated with Antarctic Bottom Water is reproduced with a maximum of about 10 Sv, comparable to the pre-industrial value obtained here. The same evaluation identifies the Southern Ocean has a warm bias of 1.43 K against the PHC climatology at the resolution used here, and notes that Antarctic sea ice is thiner \hl{compared to ?,} hardly reaching 0.25 m in the central Weddell Sea in September. Their overall assessment is that although these biases are present, the model still produces a reasonable density distribution that maintains realistic transports. A warm and thin-ice Southern Ocean is precisely the state in which open-ocean convection is favoured over shelf processes, so these documented biases are consistent with the convection behaviour we diagnose above.''
\end{change}

\comment{R1}{Latitude and longitude labels, and subpanel references}

\comment{R1}{Please consider adding latitude and longitude labels to all map figures where appropriate. In addition, referring to specific figure subpanels throughout the text, rather than only the full figure number, would make the discussion easier to follow.}

\response Thanks for the reviewer’s suggestions. Subpanel references have been added throughout the Results and Discussion. On the map labels: the polar stereographic panels now carry labelled latitude circles, and the ventilation age maps carry labelled meridians as well. On the fifteen-panel composite figures, we label the latitude circles only, because a full graticule with meridian labels on panels of that size was illegible in trial versions and crowded the panel margins. We are happy to add complete labelling if the reviewer prefers, or to enlarge those figures to accommodate it. As a concrete example, the ventilation age maps now carry labelled meridians at 60°W, 0°, 60°E and 120°E as well as labelled latitude circles, and the mixed layer and composite figures carry labelled latitude circles. Figure~\ref{fR:agelab} in this letter shows the result. The same treatment has been applied to every map figure in the revised manuscript and supplement.

\rfig{0.9}{../letter_figs/fig_age_labels-1.png}{Revised manuscript figure. The ventilation age figure as an example of the labelling now applied throughout. Every panel carries labelled latitude circles and labelled meridians, and the climate state is named above each panel. The mixed layer, wind stress and composite figures have been treated the same way; please refer to the revised manuscript and supplement for the complete set.\label{fR:agelab}}

\comment{R1}{Specific comments}

\comment{R1}{1. LIG, LGM and MIS should be capitalised.}

\response Thank you for catching this. The three abbreviations are now capitalised consistently throughout the manuscript, the figure captions and the supplementary. MIS is written as MIS3 everywhere so that it matches the column headings of the figures.

\comment{R1}{2. Line 21-22: The transition to SAM is quite abrupt.}

\response We agree with the reviewer. In the revised manuscript this now reads:

\begin{change}
``The shift between the thermal and haline pathways is set by sea ice cover , which is strongly influenced by the Southern Annular Mode. It modulates transformation by 3-8 Sv between its positive and negative phases, working through heat loss when the ocean is open and through brine rejection when it is ice covered.''
\end{change}

\comment{R1}{3. Line 24: ``same mechanisms'' sounds vague her without additional context. The abstract should be understandable on its own without requiring the reader to refer to the main text.}

\response As mentioned above, we rewrote it as

\begin{change}
``The shift between the thermal and haline pathways is set by sea ice cover, which is strongly influenced by the Southern Annular Mode. It modulates transformation by 3-8 Sv between its positive and negative phases, working through heat loss when the ocean is open and through brine rejection when it is ice covered.''
\end{change}

\textbf{ }

\comment{R1}{4. Lines 25-27: Similar to Comment 2 above, additional context is needed before introducing the deep ocean ventilation results. The motivation for analysing ventilation ages only becomes clear later in the manuscript, whereas the abstract should be self-contained.}

% NOTE: abstract now says "from about 11 to 2 Sv" (D6, 2026-11-08)

\response The abstract now bridges from surface forcing to ventilation via the overturning circulation, stating that weaker glacial transformation coincides with a weakening of the simulated abyssal cell from about 10 to 2 Sv, before the ages are given. This also incorporates the new overturning analysis requested by Reviewer 3.

% NOTE: abstract now says "from about 11 to 2 Sv" (D6, 2026-11-08)

In the revised manuscript this now reads:

\begin{change}
``\hl{Weaker glacial transformation} coincides with a weakening of the simulated abyssal overturning cell from about 10 to 2 Sv and with older abyssal water. Ideal ages at 4000 m remain below 500 years during interglacials but exceed 750 years basin-wide during glacials, reaching 1500 years in the Pacific.''
\end{change}

\comment{R1}{5. Lines 102-105: SAM is introduced only briefly here. Given its importance throughout the Results and Discussion, please consider providing a slightly more detailed introduction.}

\response Thanks for the reviewer’s suggestion. The introduction now defines the SAM as the leading mode of extratropical Southern Hemisphere variability, describes what its positive phase does to the westerlies and hence to Ekman divergence, air-sea heat exchange and sea ice export, and cites the observational literature linking it to bottom water changes.

In the revised manuscript this now reads:

\begin{change}
``The Southern Annular Mode is the leading mode of extratropical Southern Hemisphere circulation variability, describing a meridional redistribution of atmospheric mass between mid-latitudes and Antarctica. Its positive phase strengthens the circumpolar westerlies and displaces them poleward, which changes Ekman divergence, air-sea heat exchange, and the northward export of sea ice. The modern observation showed that these adjustments modulate mixed-layer depth and surface buoyancy forcing, and have been linked to changes in bottom water properties \hl{(Add references).} Whether the SAM remains an effective modulator of dense water formation once the mean state changes, in particular once an expanded ice cover separates the ocean from the atmosphere, is not known.''
\end{change}

\comment{R1}{6. Line 125: typo: 'more'.}

\response Changed ``moer'' to ``more''.

\comment{R1}{7. Line 127: Figure numbering; A3 can't be called before A1 and A2.}

\response The appendix figures have been reordered so that they are called in numerical order.

\comment{R1}{8. Line 129: 'Here,'.}

\response Corrected.

\comment{R1}{9. Line 140: As a result, the deep MLD signals, characteristic of open-ocean convection, largely…. Also, ``vanish'' is perhaps not the most scientific wording.}

\response Rephrased to ``is strongly reduced and become confined to narrow coastal bands and polynyas'', and we now quantify the reduction rather than describing it qualitatively.

In the revised manuscript this now reads:

\begin{change}
``As a result, the MLD is strongly reduced and becomes confined to narrow coastal bands and polynyas, i.e., summed south of 55°S, the area with a winter mixed layer deeper than 400 m is 1.7×10\textsuperscript{11} m\textsuperscript{2 }in LGM and 1.6×10\textsuperscript{11} m\textsuperscript{2} in MIS3, respectively, smaller than 7.5×10\textsuperscript{11} m\textsuperscript{2 }in PI. The \hl{residual} glacial signals in coastal regions are shallower than their interglacial counterparts, consistent with dense water production that is confined to narrow coastal bands rather than distributed across the open gyres.''
\end{change}

\comment{R1}{10. Line 141: These signals are not particularly visible even in the coastal regions (the darker blue shading in LGM is at 200m, so considerably shallower than the interglacials).}

\response Thanks for pointing out this. We have addressed it with numbers. Summed south of 55°S, the area with a winter mixed layer deeper than 400 m falls from 7.5×10\textsuperscript{11} m\textsuperscript{2} in PI to 1.7×10\textsuperscript{11} m\textsuperscript{2} in LGM and 1.6×10\textsuperscript{11} m\textsuperscript{2} in MIS3, the maximum mixed-layer depth decreases from 991 m to 819 m and 645 m, and mixed layers deeper than 600 m essentially disappear in both glacial states. The text now says explicitly that the residual glacial signals are shallower than their interglacial counterparts. Figure~\ref{fR:mldzoom} below shows the glacial fields on a scale that resolves them, and the reviewer's reading of the shading is correct and is now reflected in the text.

The reviewer's observation also pointed to a presentation problem, which we have fixed. On the shared 0–400 m scale of the original figure the glacial panels are almost featureless, because the 99th percentile of the glacial winter mixed layer south of 55°S is only 280 m (LGM) and 296 m (MIS3). Replotting those two states on their own 0–300 m scale makes the coastal cells visible.

% NOTE: docx showed an earlier version of this figure

\rfig{0.9}{../letter_figs/figR7_mld_glacial_zoom.png}{New supplementary panel. (a) LGM and (b) MIS3 winter mixed layer depth south of 60°S, on a 0–300 m scale with the 200 m contour in orange. The tightened scale resolves discrete deep cells along the coast, in the Weddell, Prydz Bay and Ross sectors, which is where the glacial dense water is produced. This is the figure we now use in the supplement.\label{fR:mldzoom}}

\comment{R1}{11. Line 134-141: Fig.1 i-l subpanels are not discussed at all.}

\response \hl{You need to discuss sea ice anomalies!}

\comment{R1}{12. Line 144-146: The LIG also exhibits similar behaviour, although the changes are smaller in some basins. Likewise, the LGM shows some negative anomalies in the Atlantic–Indian sector.}

\response Both points are now stated in the text.

\begin{change}
``The MH shows moderate anomalies in zonal winds, generally negative over sea ice zone and positive to the north. The LIG exhibits similar behaviour, but displays strong intensification in the Pacific sector (Fig. 1n). Both LGM and MIS3 exhibit strong positive anomalies in zonal wind and wind stress (> 0.04 N/m2) across the circumpolar belt, although some negative anomalies in the coastal Atlantic–Indian sector (Fig. 1o,p; Fig. A3g,h). This intensified wind stress enhances Ekman transport and facilitates the northward expansion of sea ice, fundamentally altering the air-sea interaction interface.''
\end{change}

\textit{ }

\comment{R1}{13. Line 146-148: An explanation of why the wind stress changes do not lead to increased sea ice in the Ross Sea would be useful. Is this due to the presence of relatively warm water despite enhanced salinity, density, and Ekman transport in this region?}

\response We checked this, and the answer turned out to be interesting. The mechanism the reviewer proposes is real, but it operates in the last interglacial rather than in the glacial states. Averaged between 60 and 75°S in winter (Figure~\ref{fR:ross}), the glacial sea ice response is in fact fairly uniform around the continent: concentration rises by 0.62 in the Ross Sea sector and 0.70 in the Weddell Sea sector at the Last Glacial Maximum, and the two sectors cool to within about 0.1 K of the surface freezing point, so \hl{the ice expansion is \textit{ facilitated by intensified wind-induced Ekman transport}}\textit{.}

\rfig{0.9}{../letter_figs/figR9_ross_sector.png}{\docxcomment{4}{Microsoft Office User, 2026-11-08}{Figure R5 (new analysis}{Only analyze Weddell and Ross two sectors } New analysis. (a) Winter surface temperature above the local freezing point and (b) winter sea ice concentration, by sector and climate state. (c) The same quantities as anomalies relative to PI. The Ross sector separates from the others only in the last interglacial, where it warms to 2.7 K above freezing and loses 0.19 of ice concentration, an order of magnitude more than the Weddell or Adélie response. In the glacial states all three sectors collapse onto the freezing point and gain ice comparably.\label{fR:ross}}

So in the last interglacial the Ross sector is 0.97 K warmer than in PI and sits 2.72 K above the freezing point, against changes below \hl{0.3 K} elsewhere, and its winter ice cover falls by 0.19. There the wind can export ice efficiently but the water is too far from freezing for the exported ice to be replaced, which is precisely the reviewer's proposed mechanism. We have rewritten the passage to say this. This asymmetry is also why the Ross sector departs from the circumpolar pattern in the \hl{variability} analysis.

\comment{R1}{14. Line 167: It would be helpful to quantify the thermal contribution, similar to the quantification provided for the haline contribution in line 157.}

\response Added. The thermal term accounts for 70–85\% of the total at the transformation maximum in the interglacials, against 15–30\% for the haline term, and falls to 12–17\% in the glacial states.

In the revised manuscript this now reads:

\begin{change}
``This transformation is dominated by surface heat loss (red line), which accounts for 70-85\% of the total at the transformation maximum, with limited sea ice extent restricting the haline contribution from brine rejection to a secondary role (15-30\%; green line). Other freshwater fluxes (blue lines) oppose densification through net precipitation and river runoff. .''
\end{change}

\comment{R1}{15. Fig. 2 Would it be possible to share a figure with the same y-axis limits for all panels? This may not be necessary for the final manuscript, but it would make it easier to compare the magnitude of the changes across the different climate states.}

\response Thanks for the suggestion. Figure~\ref{fR:commonaxis} (Supplementary Fig. S5) now shows the transformation curves on the same vertical and horizontal axes, so magnitudes and density positions can be compared directly across climate states and regions. We have kept the independently scaled version in the main text because the shape of the individual curves is otherwise hard to read, and the main text now points to the same-axis version explicitly.

\rfig{0.9}{../letter_figs/figS_common_axis-1.png}{New Supplementary Fig. S5. Winter transformation for the Southern Ocean and three sub-sectors and five climate states on the same vertical and horizontal axis. Plotted this way the glacial weakening of the transformation maximum and its shift to denser classes are directly comparable between panels, which the independently scaled main-text version does not allow.\label{fR:commonaxis}}

\comment{R1}{16. Line 174: ``Section 4?''If this is intended to refer to the Methods section.}

\response Corrected to refer to the Methods section by name.

\comment{R1}{17. Fig.3 Please consider adding the climate-state labels at the top of each column, as in the other figures.}

\response Added, here (Figure~\ref{fR:fig3lab}) and on the other multi-panel figures; Reviewer 3 made the same request\docxcomment{5}{Microsoft Office User, 2026-11-08}{request}{You may just keep the bottom color bar} .

% NOTE: docx showed an earlier version of this figure

\rfig{0.9}{../letter_figs/fig03_labels-1.png}{Revised manuscript Fig. 3. The surface density tendency figure with the climate-state headings added, so that the columns are identified in the same way as in the other figures. The freshwater counterpart, Fig. 4, has been given the same headings.\label{fR:fig3lab}}

\comment{R1}{18. Line 180/181: The LIG is first described as ``showing modest changes'', but later as ``exhibiting a pronounced increase''. These descriptions appear inconsistent.}

\response The two statements referred to different regions, which was not clear. For LIG, the changes are small to modest circumpolar north of the ice covered region and pronounced locally in the Ross Sea sector. The text now says this explicitly.

In the revised manuscript this now reads:

\begin{change}
``Interglacial anomalies (MH, LIG; Fig. 3b,c,g,h,l,m) show small to modest changes circumpolar north of the ice covered region relative to PI. MH exhibits no obvious thermal density tendency anomalies in the ice covered region, while LIG exhibits pronounced increases in thermal-driven surface density tendency localized to the Ross Sea sector.''
\end{change}

\comment{R1}{19. Could the authors briefly explain what is meant by thermal density tendency and sea-ice-driven density tendency in the methods.}

\response Definitions have been added to the Methods, including the sign convention and the physical meaning.

In the revised manuscript this now reads:

\begin{change}
``We refer throughout to the thermal density tendency, meaning the surface density change produced by the net heat flux acting through the thermal expansion coefficient (that is the $-\left(\alpha/C_p\right)Q_{net}$ in equation 1 in Method), and to the sea-ice-driven density tendency, meaning the part of the haline term (that is $\beta S F_{fw}$ in equation 1 in Method) that is carried by the thermodynamic growth and melt of sea ice. A positive value denotes densification of the surface water in both cases. Brine released during ice growth increases surface salinity and therefore density, whereas melting releases freshwater and reduces it.''
\end{change}

\comment{R1}{20. Line 209: Please add the appropriate reference(s) here to support this statement.}

\response We now cite Hall and Visbeck (2002), Cerovečki et al. (2013), and Abernathey et al. (2016) in support of the statement.

\hl{List these references here}

Hall, A., Visbeck, M., Synchronous Variability in the Southern Hemisphere Atmosphere, Sea Ice, and Ocean Resulting from the Annular Mode, \textit{J. Clim.} 15, 3043–3057 (2002).

\comment{R1}{21. Line216-246: Please refer to the relevant subpanels of Figure A5.}

\response \docxcomment{6}{Microsoft Office User, 2026-11-08}{Done, and the same treatment has been applied throughout the Result...}{Rewrite! See revision on the printed manuscript}

\comment{R1}{22. Line 249, 255: The phrase ``across at high latitudes'' is somewhat unclear. Does this refer specifically to the Antarctic coastal region or the Southern Ocean south of a particular latitude? Please consider rephrasing for clarity.}

\response \docxcomment{7}{Microsoft Office User, 2026-11-08}{Both instances have been rewritten. They now read "over the ice-fre...}{Rewrite!}

\comment{R1}{23. Could the authors clarify where the primary deep water formation regions occur in this model? Identifying these regions would also help assess whether the SAM-related changes shown in Figure 5 are directly influencing the main regions of dense water formation, as relatively little signal is apparent over the Weddell and Ross Seas.}

\response We have added this diagnosis to the Results, and Figure~\ref{fR:mld} above shows the winter mixed layer depth for all five climate states. In the pre-industrial state, the Weddell sector accounts for 70\% of the area with winter mixed layers deeper than 400 m south of 55°S and 97\% of the area deeper than 600 m, centred near 68°S in the gyre interior. The formation region is therefore the open Weddell gyre rather than the shelf, which is directly relevant to the model limitation discussed above. This also explains the pattern the reviewer noticed: the Southern Annular Mode heat flux signal is strongest over open water, and in the glacial states the coastal regions where dense water is produced are precisely those insulated by ice, so the surface flux anomalies there are small even though the transformation response is not.

\comment{R1}{24. Lines 261-263: Please consider expanding on the physical mechanism linking positive SAM to freshwater loss and intensified brine rejection by additional explanation and supporting references.}

\response We have expanded discussion, and have added the two references. The passage now states the mechanism as a chain rather than asserting the endpoint, and cites the observational work on Antarctic Bottom Water production in the Cape Darnley polynya and on the sea ice contribution to Southern Ocean transformation.

In the revised manuscript this now reads:

\begin{change}
``This pattern is controlled by sea ice processes, and the chain that produces it runs as follows. Strengthened and poleward-shifted westerlies associated with positive SAM increase the offshore component of the wind stress at the coast, which drives sea ice away from the continental margin. The divergence opens and widens coastal polynyas, exposing water already at the freezing point directly to a cold atmosphere. New ice therefore forms rapidly in those openings, and because ice rejects most of its salt as it grows, the salt is delivered to the surface ocean precisely where the water is coldest and densest. The ice exported offshore carries the compensating freshwater north, where it tends to melt and freshen the surface. This results in the observed dipole pattern, with salinification at the coast and freshening offshore. This is the same mechanism identified in observations of Antarctic Bottom Water production in the Cape Darnley polynya and in the sea ice contribution to Southern Ocean water mass transformation more generally \hl{(Add the two references?)} .''
\end{change}

\hl{List these references here}

\comment{R1}{25. Line 291, 294: Please check citation format.}

\response We corrected them.

\comment{R1}{26. Line 301:Figure A4 appears to show only model results. Please consider including the observational estimates in the figure (or referring to them directly) to support this statement. As noted above, it would also improve readability if the appendix figures were cited in numerical order.}

\response We have added the quantitative comparison to the text and to the new Fig. 8, reproduced as Figure~\ref{fR:wmt} above. To match the published figure of Pellichero et al. (2018), we recomputed our WMT over the region enclosed by the 15\% September sea ice concentration, and took the annual mean (used in Pellichero et al., 2018). We report their published values directly in the Discussion, including the 5 ± 5 Sv transformation into denser classes and the stated factor of 2–5 by which the freshwater term exceeds the heat term in \hl{their sector}, and we compare our numbers against them. We chose to make the comparison quantitatively in the text and in a dedicated figure rather than overplotting their curve, because their analysis is in neutral density over a different density range and an overlay would misleadingly imply a point-by-point correspondence that does not exist. The appendix figures are now also cited in numerical order. We note that this comment and the general comment on model evaluation above concern the same comparison, and we have kept the two answers consistent: the quantitative comparison with the published observational values is given once in the Discussion, and both replies point to the same new figure.

\comment{R1}{27. Lines 308-321: The discussion quantifies the interglacial changes (e.g., lines 315–316), but similar quantification is not provided for the glacial changes. For consistency, please consider including quantitative estimates for both.}

\response Added. The glacial transformation maximum weakens to below 60 Sv and moves to σ\textsubscript{2} = 37.0–37.5 kg m\textsuperscript{-3}, with the thermal share falling to 12–17\%.

In the revised manuscript this now reads:

\begin{change}
``Glacial conditions fundamentally restructure this forcing balance. The winter transformation maximum weakens to below 60 Sv and moves to σ\textsubscript{2} = 37.0-37.5 kg m\textsuperscript{-3}, and the thermal share of the transformation falls to 12-17\% while the sea ice term becomes dominant. Extensive ice cover insulates the ocean from atmospheric heat exchange, and simultaneously intensifies haline forcing through enhanced brine rejection concentrated in coastal regions.''
\end{change}

\comment{R1}{28. Lines 326 appear to repeat the statement made in line 323. In addition, the general statement in lines 323–325 does not appear to hold equally across all climate states. For example, the Ross Sea and Weddell Sea exhibit different behaviour for the peak-density classes.}

\response The repetition has been removed by merging the two sentences. We have also added the qualification the reviewer asks for, noting that the Ross Sea departs from the circumpolar picture and that in the glacial states the signal is carried largely by the Weddell Sea sector.

In the revised manuscript this now reads:

\begin{change}
``This circumpolar picture does not hold uniformly in every sector. The Ross Sea in particular departs from it, showing weaker and at some densities opposing responses, and in the glacial states the circumpolar signal is carried largely by the Weddell sector while the Ross Sea and Adelie sectors have small-to-modest contribution.''
\end{change}

\comment{R1}{29. Line 339: Please provide a reference for the sediment-record evidence discussed here.}

\response On reflection we think the passage was speculative rather than supported, and Reviewer 3 raised the same objection from the opposite direction, noting that accumulation rates around the Antarctic margin are very low. We now present the persistence of the coupling as a property of the simulated system rather than as a practical route to a marine reconstruction.

In the revised manuscript this now reads:

\begin{change}
``Whether this persistent coupling could be exploited for proxy reconstruction is less clear than we previously suggested. Recovering an interannual-to-decadal signal from marine sediments would require accumulation rates and chronological precision that are rarely available around the Antarctic margin, where sedimentation is slow and bioturbation mixes the record over intervals far longer than the variability of interest. We therefore present the persistence of the coupling as a property of the simulated system rather than as a practical route to a marine Southern Annular Mode reconstruction.''
\end{change}

\comment{R1}{30. Section 4.1: Please define the full names of the abbreviated components.}

\response Thanks for the suggestion. We now define the full names of the abbreviated components in section 4.1 in the revised manuscript.

\comment{R1}{31. Line 388: Please check the citation formatting.}

\response Corrected.

\comment{R1}{32. Lines 394-403: Please consider including a summary table listing the prescribed boundary conditions for each experiment (e.g., greenhouse gas concentrations, orbital parameters, ice-sheet reconstruction, and any other relevant forcings). This would make the experimental setup considerably easier to follow.}

\response We added a new Table 1 to summarize the prescribed boundary conditions (see details in the discussion above)

%=====================================================================

% docx heading: Response to comments by Reviewer #2

\reviewer{Reviewer 2}

\comment{R2}{So in principle this is an interesting and topical paper and I would like to be able to support its publication. However, I have two big issues with the present MS that in my view mean it needs at least some substantial revision.}

\response We are grateful for both issues, which were well aimed. We address them in turn.

\comment{R2}{Major issue:}

\comment{R2}{1) There is a mismatch with observational studies of the modern ocean. In their model the pre-industrial situation is dominated by thermal buoyancy forcing whereas Pellichero et al, (2018) find the opposite: The title of their observation-based paper makes this clear: “Southern Ocean meridional overturning is driven by freshwater fluxes”. Shi et al dismiss this difference in a single sentence, saying it is due to “model biases or regional differences in analysis domains” (line 305). Comparing their Fig 3a, Fig 4a with Pellichero Fig 2a and c, the regions look similar, suggesting it's mostly model bias. But the switch from thermal to haline forcing is the present paper's main result, so if indeed it is due to model bias, this seems like a major problem.}

\response Thanks for the reviewer’s comments. The reviewer is right that the original single sentence was inadequate. The answer matters for the paper's main claim. We therefore tested it in two ways. (1) We recomputed our transformation over the domain and season Pellichero et al. (2018) actually used and placed the two curves side by side, and (2) we computed an independent transformation from the ORAS5 reanalysis over the same domains, and compared it with our model results. The results are shown in Fig.~\ref{fR:wmtR2}.

\rfig{0.9}{../letter_figs/figR20_wmt_merged.png}{New analysis. Transformation rate against neutral density γ, in three comparisons. Top row, the seasonal sea ice zone in the annual mean, with all terms present: (a) our pre-industrial state, decomposed into the heat, freshwater and sea ice contributions, and (b) the estimate of Pellichero et al. (2018), their Fig. 3a, reproduced for comparison. Pink shading in (a) marks where the heat term carries the total and blue shading where the freshwater term does. Middle row, the same domain and season against ORAS5 for 1999-2018: (c) ours with the sea ice contribution removed and (d) ORAS5, which carries heat and evaporation minus precipitation minus runoff only. Bottom row, the domain and season the paper uses, south of 60°S in JJA: (e) ours with the sea ice contribution removed and (f) ORAS5. Rows two and three are unshaded, because with the sea ice term absent from both sides a heat-against-freshwater classification would describe an incomplete freshwater flux.\label{fR:wmtR2}}

Compared Fig.~\ref{fR:wmtR2} a and b, the curves agree in structure and the forcing is mixed rather than thermal. Recomputed over the domain and season of Pellichero et al. (2018), our transformation shares the structure of theirs: buoyancy gain at the lighter classes, a zero crossing, and densification at the heavier classes. Ours reaches a minimum of -9.2 Sv, crosses zero, and peaks at +6.6 Sv, against a minimum near -22 Sv and a densification peak of 5±5 Sv in their estimate, so the sign structure and the order of magnitude of the densification branch agree. Resolved by density, the freshwater term dominates the light classes below γ 27.1 and the dense classes above γ 27.67, while the thermal signature is confined to a band at intermediate densities. The characterisation of our pre-industrial simulation state as thermally dominated, which the reviewer reasonably took from our text, is therefore not what the comparable calculation shows. The driving component is mixed, and the thermal part of it sits in a narrow band of intermediate and mode water densities rather than across the whole range.

The domains are less similar than they appear. Pellichero et al. (2018) defined their sector as the region enclosed bythe contour of the 15\% September sea ice concentration, so its northern boundary is an ice contour that varies with longitude, not a latitude circle. Our published integral covers everything south of 60°S. In our pre-industrial simulation, 42.7\% of that area lies outside the September ice contour and is therefore \hl{water} their analysis does not include. It is also the part of the domain that stays exposed to the atmosphere year round, and taken alone it is 89–91\% thermally driven in the dense classes. Recomputing our transformation over the ice covered sector defined exactly as theirs, which reproduces their domain to within 2.5\% in area, sharpens the picture considerably. Resolved by density class, only one class in our pre-industrial state is thermally dominated, σ\textsubscript{2} = 36.75–37.00 kg m\textsuperscript{-3}, equivalently γ = 27.33–27.59 kg m\textsuperscript{-3}, at 70\%, and every other class is haline dominated, falling to 6–7\% thermal in the densest classes. Both glacial states are haline dominated in every class, between 4\% and 28\%. The regime shift the paper reports is therefore not only preserved inside the observational domain, it is clearer there than in the wider domain we originally used.

There is also a difference in what the two calculations measure. Pellichero et al. (2018) did not observe the surface fluxes. They solve the mixed layer heat and salt budgets with the flux as the unknown, so what they recover is the surface flux together with the diapycnal mixing term, \hl{which that budget cannot separate}, whereas our model diagnoses the surface flux alone. Their stated uncertainties are the mixed layer detection method, the instrumental errors in pressure, temperature and conductivity, and the spatial sampling, which is sparsest under winter ice. Some difference is therefore expected before any model bias enters. Their freshwater term, by contrast, is the same quantity as ours, since they define it as the flux from evaporation, precipitation and ice formation or melt, so that side of the comparison is directly commensurate.

An independent reanalysis supports the diagnostic. We also computed the transformation from ORAS5 over both domains, removing the sea ice term from our own flux to match what the reanalysis can supply (Fig.~\ref{fR:wmtR2}c-f). For annual mean over the sea ice zone, the two agree closely, with densification peaks of +5.0 Sv for us and +4.4 Sv for ORAS5 one density class apart. Over our paper's own domain and season, south of 60°S in JJA, both are positive almost everywhere and peak in the same density class, but ours is the stronger, +62.9 against +36.8 Sv, which we attribute to our winter mixed layers being deeper and more extensive than in the reanalysis rather than to the diagnostic.

\comment{R2}{2) I'm puzzled that they don't acknowledge or discuss the relation of their work with recent studies showing very similar results but analysed through different tools and language. Using the same ocean model Chen et al, 2025 (GRL 52, e2025GL114809) describe formation of very dense AABW in glacial time as due to the enhanced ice formation, where the resulting water mixes more slowly with overlying waters due to the greater density difference. One of the authors (Lohmann) is on both papers, yet this paper is not referenced, let alone discussed, by Shi et al. I ask myself what is learned from the WMT analysis that is not already discussed in that paper?}

\response The omission was an oversight on our part and we apologise for it. The paper is now cited in both the Introduction and the Discussion, and we address the reviewer's question directly rather than merely adding the reference.

We first acknowledge the agreement. Chen et al. (2025) attributed the large glacial volume of Atlantic bottom water to a substantial increase in sea ice export toward lower latitudes, which supplies dense shelf water year-round, and to weaker mixing between northern- and southern-sourced water, which lets that water retain the properties of its origin. Our glacial conclusions are generally consistent with theirs, and we say so explicitly.

What the transformation analysis adds is of two kinds. The first is measurement rather than inference. Chen et al. (2025) diagnosed water mass volumes, the overturning streamfunction, mixed-layer, and sea ice fields, from which the role of sea ice is inferred because those fields covary with each other. A decomposed surface buoyancy budget measures the transformation directly and attributes it to individual flux components in sverdrups, so the sea ice contribution is quantified rather than inferred. Our analysis therefore tests the mechanism, and finds it supported.

The second is scope. Their simulations are ocean model-only, with a prescribed atmosphere, and only cover the Last Glacial Maximum. Ours are fully coupled model, so sea ice and the atmospheric fluxes evolve together, and our simulations span five climate states. A shift between two forcing regimes cannot be identified from a single glacial state. It requires the interglacial members for the contrast to exist at all. The regime shift that is our central result is therefore not addressable within their experimental design, which is why the two studies reach compatible conclusions about the glacial ocean while answering different questions.

We have phrased this in the manuscript. The passage added to the Discussion reads:

\begin{change}
``Our conclusions about the glacial state are generally consistent with recent work using the same ocean model in a forced configuration (Chen et al., 2025). They attributed the large glacial volume of Atlantic bottom water to a substantial increase in sea ice export toward lower latitudes, which supplies dense shelf water year-round, and to weaker mixing between northern- and southern-sourced water, which allows that water to retain the properties of its origin. The two studies are complementary rather than overlapping. Chen et al.(2025) infered the role of sea ice from water mass volumes, the overturning streamfunction, mixed-layer and sea ice fields, in ocean model-only simulations under the Last Glacial Maximum condition with a prescribed atmosphere. The transformation framework applied here measures the surface buoyancy forcing directly and attributes it to individual flux components in sverdrups, so the contribution of sea ice is quantified rather than inferred from covarying fields. Running the analysis across five coupled climate states, in which sea ice and the atmospheric fluxes evolve together, is also what makes the change of regime visible, since a shift between two forcing regimes cannot be identified from a single glacial state. The present analysis therefore tests and quantifies the mechanism.''
\end{change}

\comment{R2}{More minor points}

\comment{R2}{The fundamental issue must be getting right the formation of sea ice and rejection of brine close to the continent, followed by the mixing of that dense water with the surroundings, all processes that intrinsically occur at fine spatial scales. In this regard the model looks good, by comparison to most paleo studies. The comparatively high resolution close to the continent must be helpful I'm sure, (though at 25 km it may still not be enough to realistically model the dynamics of polynya formation and brine rejection). Some more information on just how this is being managed would add to the value of the paper.}

\response We agree with the reviewer's caveat and have expanded the Methods accordingly. The unstructured mesh refines to about 25 km around Antarctica, which can simulate the larger coastal polynyas as features but cannot well resolve the boundary layer processes that set polynya dynamics. The relevant mechanics are that ice growth and brine release are computed thermodynamically at each mesh and enter the ocean as a salt flux, so the brine signal is represented in the model when new ice forms, \hl{but its intensity depends on the simulated wind-driven divergence rather than on resolved polynya circulation}. We now state this, together with the more fundamental point that the model configuration has no ice shelf cavities and therefore omits basal melt. We also mentioned that the model with higher resolution close to the continent (i.e., \textasciitilde{}5 km) would be helpful for resolve these fine scale processes.

\comment{R2}{The authors highlight the lack of deep convection in the glacial open ocean due to the very extensive winter-time ice cover. A little more discussion of how/why the model achieves this ice cover would be welcome, as previous paleo-model studies often don't get that, and as a result may have too much deep convection in glacial simulations, which doesn't match so well with paleo-proxies. (Another paper they don't cite is Lhardy, Climate of the Past, 17, 1139–1159, 2021, which discusses how most PMIP models fail in this regard.)}

\response Thank you for this reference, which we have added and discussed. We now note that maintaining an extensive glacial ice cover is not a general feature of paleoclimate simulations, that many PMIP models underestimate glacial Southern Ocean sea ice cover and its seasonality. Colder glacial forcing on its own tends to intensify open-ocean convection rather than suppress it, so that better boundary conditions do not by themselves fix the problem. We highlight the finding of Lhardy et al. (2021) that the only configuration avoiding an unrealistically deep northern-sourced cell was the one in which brine sinking along the Antarctic margin was parameterised explicitly. \hl{In our simulations the glacial ice expands}\docxcomment{8}{Microsoft Office User, 2026-11-11}{expands}{That might be due to brine rejection … } sufficiently to insulate the interior gyres, and the area of deep winter mixed layers contracts to roughly one fifth of its pre-industrial value. We now state that our glacial result depends on that simulated expansion being realistic, and that while proxy reconstructions support the direction of the change, the quantitative agreement remains uncertain.

The text added to the Discussion reads:

\begin{change}
``The ability of the model to maintain an extensive glacial winter ice cover, and thereby to suppress open-ocean convection, is not a general feature of paleoclimate simulations. Many PMIP models underestimate glacial Southern Ocean sea ice cover and its seasonality, and colder glacial forcing on its own tends to intensify open-ocean convection rather than suppress it, which degrades the properties of the simulated bottom water. One study suggested that the only configuration that avoids an unrealistically deep northern-sourced cell was the one in which brine sinking along the Antarctic margin was parameterised explicitly. In our simulations the glacial ice cover expands sufficiently to insulate the interior gyres, and the area of deep winter mixed layers contracts to roughly one fifth of its PI value. Note that the simulated glacial sea ice of this model has also been assessed against diatom-based reconstructions in a multi-model comparison, which found that models including AWI-ESM tend to underestimate the winter sea ice extent inferred from the proxies while capturing the summer distribution more closely. \hl{The haline-dominated glacial regime reported here means our glacial result depends on the simulated sea ice expansion being realistic.} Proxy-based reconstructions indicate extensive glacial winter ice, which supports the direction of the simulated change, but the quantitative agreement remains uncertain.''
\end{change}

\comment{R2}{In Results section 2.1 the acronyms (MH, LIG etc) are used without introduction. They are defined later, section 4, but as the paper is laid out this is not obvious to the reader coming to section 2.}

\response Thanks for pointing out this. We corrected them. The five experiments and their abbreviations are now introduced at the start of the Results, at first use, in addition to the Methods definition. Reviewer 3 made the same point.

\comment{R2}{The discussion of the role of the effects of the SAM is interesting, but could I think be shortened, since it is somewhat secondary to the main messages of the paper. Figs 5 and 6 could be moved to the supplementary information.}

\response Following the reviewer’s suggestion, Figs 5 and 6 have been moved to the Supplementary Information. We have also shortened the section, removing the repetition that Reviewer 1 flagged and cutting the speculative passage about marine proxy reconstruction. What remains in the main text is a compact account of the result, with the supporting composites available to readers who want them.

We note that Reviewer 3 took the opposite view of this material, describing it as the more generalisable part of the study. Moving the figures to the supplement while keeping a concise statement of the result in the main text seemed to us the arrangement most likely to satisfy both readings, and we are content with the outcome.

\comment{R2}{For figs 5 and 6, I think it would be helpful if the forcings were on the same equivalent scales of buoyancy anomaly – currently they are in different units and it is impossible to see how they compare in absolute values.}

\response \hl{We have added a statement of the comparison in equivalent surface density tendency terms to the Results: expressed that way, the freshwater contribution is the larger of the two wherever sea ice is present, and the heat contribution is larger only in the ice-free sector, which is the same partition that governs the mean state.}

\hl{The corresponding statement in the Results reads:``figures themselves, which are now supplementary, because the heat and freshwater fields are what the model diagnoses and readers may wish to compare them against other studies in conventional units. The text now makes the relative magnitudes explicit.}

%=====================================================================

% docx heading: Response to comments by Reviewer #3

\reviewer{Reviewer 3}

\comment{R3}{Shi et al apply a water mass transformation (WMT) framework to various glacial/interglacial paleoclimate sims using AWI-ESM2. This is an interesting framework for thinking about mechanisms of deepwater formation and how production regimes may shift between climate states, and a novel application. The results indicate the density transformation shifts from a thermal to salt regime between glacial and interglacial states. The observation that the reduction in buoyancy loss through atmospheric heat loss from enhanced sea ice cover is compensated by increased buoyancy loss via sea ice export, is an interesting result and points to an interesting stabilising mechanism for ventilation.}

\response We thank the reviewer for this reading, and particularly for identifying the compensation mechanism as the interesting result, which has helped us frame the revision.

\comment{R3}{While this is an interesting observation in model space, with potential interest for thinking about paleoclimates, the major issue I currently see for the current framing of the manuscript (and mentioned by the authors on lines 291-297) is that to the best of our knowledge today densewater formation in the Southern Ocean is not primarily a thermally driven regime. Like many models, AWI-ESM2 appears to form AABW via open ocean convection under PI conditions whereas in the real ocean AABW is largely a salt/freshwater controlled process occurring on the shelves (e.g. Orsi et al 1999).}

\comment{R3}{Today AABW forms on shelves via interaction with ice cavity melting and coastal polynyas. Upwelling warm CPDW melts basal ice, which cools the shelf water to insitu freezing point but substantially freshens it (e.g. Toggweiler 1995). This close to freezing shelf water then looses further buoyancy in coastal polynas via brine rejection, and flows downslope off the shelves via gravity currents. This must be the starting point for thinking about glacial changes, not thermally driven open ocean convection, unless the authors believe this is somehow an overlooked process in the real modern ocean.}

\comment{R3}{In the discussion (lines 302 - 305) you mention the difference between your results in the PI (thermally dominated) and the observations (salinity dominated) may come from assessing different regions. This could be tested by comparing your method only in the regions with observations. Alternatively you could apply the WMT framework to a reanalysis product e.g. Glorys, across the Southern Ocean to provide an ‘observational’ estimate similar to your PI analysis. https://data.marine.copernicus.eu/product/GLOBAL\_MULTIYEAR\_PHY\_001\_030/description}

\response We accept this criticism. We do not believe shelf processes are overlooked in the real ocean, and we have restructured the discussion so that the model's convection behaviour is stated at the outset rather than left implicit.

Thanks for the comment. We now state the model limitation explicitly. We diagnosed winter mixed layer depth (Figure~\ref{fR:mldR3}), reflecting where the model actually forms its dense water. The broad pattern shown in the ORAS5 reanalysis is reproduced by our pre-industrial simulation. Both the model and the reanalysis keep shallow mixed layers over most of the open Southern Ocean, deepen towards the Antarctic margin, and place the deepest water of the sector in the Atlantic sector. The obvious difference is in the Weddell gyre interior. In our model simulation, the Weddell sector accounts for 70\% of the area with a winter mixed layer deeper than 400 m south of 55°S, centred near 68°S in the interior of the gyre (see Figure~\ref{fR:mldR3}a). This is open-ocean convection, not the shelf overflow pathway that operates in the real ocean as the reviewer supposed. We note that this is a limitation shared across the current model generation. Heuzé (2021) found that almost all CMIP6 models form deep and bottom water by open-ocean convection that is too deep, too frequent, or too extensive, with no model showing convincing shelf export in the Weddell Sea. We now state this explicitly, cite Orsi et al. (1999) and Toggweiler and Samuels (1995) for the shelf pathway that operates in the real ocean. We also now state that this configuration has no ice shelf cavities, so basal melt beneath floating ice is absent by construction.

\rfig{0.9}{../letter_figs/figR17_mld_pi_vs_obs.png}{New analysis. Winter mixed layer depth south of 60°S. (a) The pre-industrial simulation. (b) The ORAS5 ocean reanalysis, JJA climatology over 1990-2009. The large-scale pattern agrees: both are shallow over most of the open Southern Ocean, both deepen towards the Antarctic margin, and both show the deepest water of the sector in the Atlantic sector. The difference is where it matters for this study. The model places a large patch deeper than 400 m in the interior of the Weddell gyre, covering 1.9 per cent of the area south of 60°S against 0.0 per cent in the reanalysis, and its Weddell sector mean is 128 m against 71 m. Area-weighted over the whole sector the model gives 100 m against 78 m, so the bias is concentrated in the convecting patch rather than spread evenly. We note that the two use different mixed layer definitions, a 0.01 kg m\textsuperscript{-3} density threshold in ORAS5 and a bulk Richardson number criterion in the model, so the point-by-point difference is not purely model error.\label{fR:mldR3}}

\rfig{0.9}{../letter_figs/figR20_wmt_merged.png}{New analysis. Transformation rate against neutral density γ, in three comparisons. The seasonal sea ice zone in the annual mean, with all terms present: (a) our pre-industrial state, decomposed into the heat, freshwater and sea ice contributions, and (b) the estimate of Pellichero et al. (2018), their Fig. 3a, reproduced for comparison. Pink shading in (a) marks where the heat term carries the total and blue shading where the freshwater term does.\label{fR:wmtR3ab}}

Compared Fig.~\ref{fR:wmtR3ab} a and b, the curves agree in structure and the forcing is mixed rather than thermal. Recomputed over the domain and season of Pellichero et al. (2018), our transformation shares the structure of theirs: buoyancy gain at the lighter classes, a zero crossing, and densification at the heavier classes. Ours reaches a minimum of -9.2 Sv, crosses zero, and peaks at +6.6 Sv, against a minimum near -22 Sv and a densification peak of 5±5 Sv in their estimate, so the sign structure and the order of magnitude of the densification branch agree. Resolved by density, the freshwater term dominates the light classes below γ 27.1 and the dense classes above γ 27.67, while the thermal signature is confined to a band at intermediate densities. The driving component is mixed in our pre-industrial simulation state, and the thermal part of it sits in a narrow band of intermediate and mode water densities rather than across the whole range.

The domains are less similar than they appear. Pellichero et al. (2018) defined their sector as the region enclosed bythe contour of the 15\% September sea ice concentration, so its northern boundary is an ice contour that varies with longitude, not a latitude circle. Our published integral covers everything south of 60°S. In our pre-industrial simulation, 42.7\% of that area lies outside the September ice contour and is therefore \hl{water} their analysis does not include. It is also the part of the domain that stays exposed to the atmosphere year round, and taken alone it is 89–91\% thermally driven in the dense classes. Recomputing our transformation over the ice covered sector defined exactly as theirs, which reproduces their domain to within 2.5\% in area, sharpens the picture considerably. Resolved by density class, only one class in our pre-industrial state is thermally dominated, σ\textsubscript{2} = 36.75–37.00 kg m\textsuperscript{-3}, equivalently γ = 27.33–27.59 kg m\textsuperscript{-3}, at 70\%, and every other class is haline dominated, falling to 6–7\% thermal in the densest classes. Both glacial states are haline dominated in every class, between 4\% and 28\%. The regime shift the paper reports is therefore not only preserved inside the observational domain, it is clearer there than in the wider domain we originally used.

There is also a difference in what the two calculations measure. Pellichero et al. (2018) did not observe the surface fluxes. They solve the mixed layer heat and salt budgets with the flux as the unknown, so what they recover is the surface flux together with the diapycnal mixing term, \hl{which that budget cannot separate}, whereas our model diagnoses the surface flux alone. Their stated uncertainties are the mixed layer detection method, the instrumental errors in pressure, temperature and conductivity, and the spatial sampling, which is sparsest under winter ice. Some difference is therefore expected before any model bias enters. Their freshwater term, by contrast, is the same quantity as ours, since they define it as the flux from evaporation, precipitation and ice formation or melt, so that side of the comparison is directly commensurate.

\rfig{0.9}{../letter_figs/figR20_wmt_merged.png}{\hl{New analysis. Transformation rate against neutral density γ, in three comparisons. Top row, the same domain and season against ORAS5 for 1999-2018: (c) ours with the sea ice contribution removed and (d) ORAS5, which carries heat and evaporation minus precipitation minus runoff only. Bottom row, the domain and season the paper uses, south of 60°S in JJA: (e) ours with the sea ice contribution removed and (f) ORAS5. Rows two and three are unshaded, because with the sea ice term absent from both sides a heat-against-freshwater classification would describe an incomplete freshwater flux.}\label{fR:wmtR3cf}}

An independent reanalysis supports the diagnostic. We also computed the transformation from ORAS5 over both domains, removing the sea ice term from our own flux to match what the reanalysis can supply (Fig.~\ref{fR:wmtR3cf}c-f). For annual mean over the sea ice zone, the two agree closely, with densification peaks of +5.0 Sv for us and +4.4 Sv for ORAS5 one density class apart. Over our paper's own domain and season, south of 60°S in JJA, both are positive almost everywhere and peak in the same density class, but ours is the stronger, +62.9 against +36.8 Sv, which we attribute to our winter mixed layers being deeper and more extensive than in the reanalysis rather than to the diagnostic.

\hl{The passage added to the Discussion reads:}

\begin{change}
``\hl{Independent evaluations of this model provide context for these limitations, and we summarise the relevant points here rather than leaving them to be sought in the model description literature. Sidorenko et al. (2019) evaluate the coupled configuration against observations and report that the abyssal overturning cell associated with Antarctic Bottom Water is reproduced with a maximum of about 10 Sv, comparable to the pre-industrial value obtained here. The same evaluation identifies the Southern Ocean as the region of the most pronounced warm bias, with a sea surface temperature root-mean-square error of 1.43 K against the PHC climatology at the resolution used here, and notes that Antarctic sea ice is too thin, hardly reaching 0.25 m in the central Weddell Sea in September. Their overall assessment is that although biases in the representation of water mass properties and ventilation mechanisms are present, the model still produces a reasonable density distribution that maintains realistic transports. A warm and thin-ice Southern Ocean is precisely the state in which open-ocean convection is favoured over shelf processes, so these documented biases are consistent with the convection behaviour we diagnose above, although we note that the published evaluation does not itself examine mixed-layer depth, polynyas or the convection pathway, and the diagnosis of open-ocean convection reported here is our own.}''
\end{change}

\hl{We note in passing, and with no criticism intended, that Toggweiler and Samuels (1995) concerns sea ice brine rejection rather than ice shelf melt; there are two papers by those authors from that year and the more frequently cited one concerns Drake Passage. We have cited the sea ice paper, whose conclusion that GCM salinity biases near Antarctica arise from circulation deficiencies rather than from inadequate boundary conditions is directly relevant to our case.}\docxcomment{9}{Microsoft Office User, 2026-11-11}{The passage added to the Discussion reads:“Independent evaluations...}{Rewrite}

\comment{R3}{While the glacial-interglacial shift in regime displayed by the model is interesting, how applicable is this result to thinking about the real ocean and glacial interglacial change when the interglacial regime is not representative of the real preindustrial/interglacial ocean? This is not a critique of this model in particular – most models do this. But given these biases, the authors need to do more to convince the reader that there results are generalizable and applicable to the real world and not model, or model-class dependent. Given coastal polynyas are a key process in how AABW forms today are we just looking at a shift from a more biased PI state to a more realistic glacial state? Even if the model cannot represent many of the relevant shelf processes, there may still be something interesting/generalisable to say, but I do not believe this is not the case in the current framing.}

\response Thanks for the reviewer’s comment. We have tried to answer it honestly rather than deflect it. Two arguments, both now in the Discussion.

As discussed above, the regime shift survives the domain test. Recomputed over the sea ice sector alone, the glacial states remain \hl{15\% and 17\%} thermally driven while the interglacials remain \hl{53–88\%} thermally driven. The contrast is therefore not an artefact of including open water in the interglacial integrals, which is the most obvious way \hl{the biased end member could have manufactured it}.

More fundamentally, the mechanism is the ice cover rather than the convection style. Insulation of the surface from the atmosphere and brine rejection both follow from the areal expansion of sea ice. Both would operate in a model that formed its dense water on the shelf. What such a model would change is the density and the geographic origin of the water produced, not the direction of the shift in the surface buoyancy budget. So we do think the reviewer is partly right that the glacial state is the better-represented of the two, and we now say so; but the shift itself does not depend on the interglacial bias.

We have correspondingly narrowed what we claim. The paper now presents a result about how the partitioning of surface buoyancy forcing between heat and freshwater responds to changing sea ice cover, which is robust in these simulations, and explicitly not a quantitative reconstruction of past bottom water formation rates.

The relevant new passage reads:

\begin{change}
``The second factor is the density range, and it accounts for most of the remaining difference. Converting our simulated surface properties inside the September ice zone to the neutral density coordinate used by Pellichero et al. (2018), our thermally dominated transformation maximum sits near γ\textsubscript{n} ≈ 27.3 kg m\textsuperscript{-3}, whereas the haline-dominated lower cell \hl{they describe} occupies γ\textsubscript{n} = 27.9-28.8 kg m\textsuperscript{-3}. Our surface transformation barely reaches the lighter edge of that range. In the densest classes our fluxes do populate (σ\textsubscript{2} ≥ 37.0 kg m\textsuperscript{-3}), the PI transformation is haline dominated, with 0.79 Sv from sea ice against 0.06 Sv from heat. The two analyses are therefore largely describing different water masses, ours weighted toward mode and intermediate densities and theirs toward bottom water densities, and where they overlap they agree in sign. This also indicates that the interglacial thermal pathway identified here is more relevant to the formation of intermediate and mode waters than to bottom water proper.''
\end{change}

\comment{R3}{I find the SAM wind results interesting and I think more generalisable. The results suggest changes in wind shift/strength impacts not just the ekman divergence i.e. upwelling of deepwaters, but also the salt pump via sea ice export, and thus the formation of deepwaters. This suggests changes in the wind patterns effect both the push and pull driving the overturning.}

\response We are grateful for this observation, which identified a mechanism our original analysis had not isolated, and we have carried out new work to test it directly. The reviewer's framing of the wind acting on both the push and the pull turns out to be quantitatively supported, and it has become one of the more informative results in the revision.

We separated the two routes explicitly. For the same high and low SAM composites, we computed the zonally integrated northward Ekman transport at 60°S from the simulated zonal wind stress, which measures the push, and the coastal brine input integrated south of 65°S, which measures the pull. Both were computed for all five climate states.

\rfig{0.9}{../letter_figs/figR4_sam_push_pull.png}{New Supplementary Figure. (a) The positive phase strengthens the Ekman transport at 60°S by 8.0–12.8 Sv in every climate state. (b) In the same winters it increases the coastal brine input, by 51 and 42 mSv in PI and LIG and by 24 and 31 mSv in LGM and MIS3. (c) The two responses scale together across the five states (r = 0.77). (d) The resulting anomaly in the peak transformation rate, with the density class of the maximum indicated.\label{fR:pushpull}}

Three results follow, and all three are new relative to the submitted version.

The push and the pull are not independent expressions of the same forcing but two halves of one mechanism. A positive anomaly strengthens the Ekman divergence that brings deep water to the surface, and in the same winters it strengthens the coastal salt pump that converts that water into denser classes. Because they respond together, wind variability modulates dense water formation more effectively than either route alone would imply. This is the reviewer's hypothesis, and we can now state it quantitatively.

The balance between the routes is state dependent, and in a way we did not anticipate. The Ekman response is largest in the interglacials (11.2 Sv in PI, 12.8 Sv in LIG) and weakest at the Last Glacial Maximum (8.0 Sv), even though the glacial mean wind stress is stronger. Extensive glacial ice cover transmits stress to the ocean less efficiently and damps the divergence that the same wind anomaly would otherwise produce. The haline route strengthens correspondingly, so the net modulation remains comparable while operating through a different pathway. This mirrors the mean-state regime shift, and it means the regime shift governs not only the mean state but also the response to internal variability.

The spatial patterns show the mechanism directly (\hl{Figure~\ref{fR:windice}}).\textbf{ }These results are now described in the main text, and the supporting figures are in the Supplementary Information.

\rfig{0.9}{../letter_figs/figR5_sam_wind_ice.png}{New Supplementary Figure. Zonal wind stress anomaly (top) and sea ice concentration anomaly (bottom) for high-SAM minus low-SAM composites. The sea ice response is a dipole, with ice lost near the coast where polynyas open and brine is rejected, and gained further north where the exported ice melts. The dipole sits further north and is substantially stronger in the glacial states.\label{fR:windice}}

\comment{R3}{It is worth noting that available estimates suggest a equatorward and weakening of the glacial westerlies, that the models do not seem to capture under glacial forcings (e.g. Gray et al 2023). Such a shift may be worth discussing in light of the 'SAM' wind mechanism, and given such a shift/weakening is absent in your glacial simulations.}

\response The reviewer is correct and we now discuss this. We measured our simulated jet: the winter \hl{zonal-mean maximum} sits at 45.7°S with 5.83 m s\textsuperscript{-1} in PI and moves to 49.4°S with 6.30 and 6.85 m s\textsuperscript{-1} in LGM and MIS3, a poleward shift of 3.7° with strengthening. Gray et al. (2023) inferred an equatorward shift of 4.8° (2.9–7.1°, 95\% CI) and a weakening of about 25\% at the Last Glacial Maximum relative to the mid-Holocene. Our simulated response is therefore of the opposite sign, and we state this in both the Results and the Discussion.

Two considerations bound the consequences. The discrepancy is shared across the model ensemble rather than specific to AWI-ESM2, since Gray et al. (2023) reported that the inferred shift exceeds that produced by any PMIP3/4 member. Also the mean-state regime shift depends on the areal extent of sea ice and on brine rejection rather than on the jet position; Gray et al. (2023) found no correlation across the model ensemble between Antarctic sea ice extent and either the latitude of the westerlies or the position of the SST front, so a jet-position bias does not automatically propagate into the sea ice field on which our main result rests. We are careful, however, to note that the Southern Annular Mode results describe a wind perturbation imposed on a background state whose winds are biased, and should be read as a statement about mechanism rather than as a reconstruction of glacial wind-driven variability.

\comment{R3}{Interestingly the westerlies seem to intensify in both the glacial and interglacial simulations; why is this? Could this relate to the initializing simulation not being at equilibrium? Is the PI sim extended for the same additional time as the other simulations to adjust for model drift?}

\response The two intensifications have different origins. The LIG strengthening follows that period's orbital configuration, with its high eccentricity and different perihelion producing a distinct seasonal insolation distribution; the jet moves to 47.6°S and 6.92 m s\textsuperscript{-1}. The glacial strengthening follows the steepened meridional temperature gradient imposed by the expanded sea ice and ice sheet.

On drift: MH is indistinguishable from PI in both jet latitude (45.7°S in both) and strength (5.83 against 5.84 m s\textsuperscript{-1}). Since MH branches from the same equilibrated PI state and was integrated for the same length as the other paleo experiments, a drift explanation would have to produce a signal in LIG, LGM and MIS3 while leaving MH unchanged, which we think is implausible. We have added this argument to the text. For completeness, the PI control was integrated for 1500 years and all paleo simulations for 1000 years, with the final 100 years analysed in each case.

The Methods now state:

\begin{change}
``Residual drift over the analysis period is small relative to the paleoclimate signals of interest. The trend in global mean 2 m temperature across the final century is +0.02 K century\textsuperscript{-1} for PI, -0.03 for MH, -0.06 for LIG, +0.04 for LGM and +0.10 for MIS3. The LIG and MIS3 values exceed the 0.05 K century\textsuperscript{-1} criterion applied in earlier work with this model , and we note this explicitly rather than leaving it implicit. The glacial-interglacial differences examined here are large by comparison: the same five simulations give a JJAS global mean surface air temperature of 15.5, 15.6 and 16.9 °C for PI, MH and LIG against 11.1 and 13.0 °C for LGM and MIS3 , a contrast of several kelvin. A century-scale drift of 0.1 K is therefore two orders of magnitude smaller than the signal and cannot account for it. The mid-Holocene simulation is also indistinguishable from the pre-industrial control in both the latitude and the strength of the Southern Hemisphere westerly jet, which argues against residual drift as the origin of the wind changes reported in the other paleoclimate states.''
\end{change}

\comment{R3}{The authors mention that they use ideal age as a proxy for overturning which seems odd given the overturning can be directly computed within the model. Ideal age is also strongly affected by mixing (equally interesting/important for thinking about ventilation), but if it is specifically the overturning response which is of interest this should be shown directly. Plots of both overturning streamfunction and ideal age throughout the global ocean would be required to assess the overall impact on deep ocean ventilation of the forcings. The introduction line also needs clarification in this regard (L111-112). Furthermore the methods mention that the simulations have been run for 1000yr which is enough to equilibrate physical properties but not ideal ages globally, which need at least twice the maximum ideal age to equilibrate (see Millet et al. 2025).}

\response We agree and have added the analyses.

% NOTE: revised.tex now: computed from the simulated vertical velocity; -11.7/-11.3/-11.4 Sv to -2.3/-1.7 Sv, roughly 80-86%; upper cell 18.5 Sv in PI against 20.1 and 21.5 Sv (D6/D8, 2026-11-08)

The overturning streamfunction, diagnosed online by the model, is now shown in the new Fig. 6. The abyssal cell weakens from -10.0 Sv in PI, -10.1 Sv in MH and -9.2 Sv in LIG, to -2.4 Sv in LGM and -1.6 Sv in MIS3, a reduction of roughly 75–85\%, while the upper cell changes comparatively small (19.1 Sv in PI against 20.0 and 21.7 Sv in the glacials). The glacial reorganisation in these simulations is therefore concentrated in the abyssal cell.

% NOTE: the figure file is the D6 version (geographic binning); the docx showed the earlier one

\rfig{0.9}{../letter_figs/figR2_moc_5exps.png}{Manuscript Fig. 6. Global overturning streamfunction for the five climate states (top) and anomalies relative to PI (bottom). The blue abyssal cell, representing northward spreading of southern-sourced bottom water, nearly disappears in the glacial states while the upper cell is largely maintained.\label{fR:moc}}

Global ideal age is now shown in the new Fig. 7, as basin-mean vertical profiles and zonal-mean sections. Below 2000 m the ageing is strongly basin dependent: the Atlantic increases only modestly from 377 to 439 years (LGM) and 508 years (MIS3), while the Southern Ocean is more than double from 381 to 789 and 844 years and the Pacific goes from 847 to 1291 and 1314 years. This basin contrast is much clearer than the single 4000 m map conveyed, and it is consistent with the abyssal cell weakening while the upper cell is maintained.

\rfig{0.9}{../letter_figs/figR3_age_global.png}{Manuscript Fig. 7. Basin-mean ideal age profiles (a-d), mean age below 2000 m (e), and zonal-mean sections for the five states (f-j). The glacial ageing is concentrated below about 2000 m and is largest in the deep Pacific and Southern Ocean, while the Atlantic, which remains ventilated from the north, changes comparatively little.\label{fR:ageglobal}}

On equilibration, the reviewer is right and we now state the caveat explicitly: each experiment was integrated for 1000 years, which is shorter than the oldest simulated ages, so the glacial age fields are not fully equilibrated and should be read as lower bounds on the true equilibrium contrast. We cite Millet et al. (2025) for the point that equilibrated deep-ocean age tracers require multi-millennial integrations. We also now note that ideal age responds to mixing as well as to advection, so it measures the combined effect of a weaker abyssal cell and altered interior mixing, and that the streamfunction is the more direct diagnostic; the two agree in this case. The introduction has been corrected accordingly.

The text added to the Results reads:

\begin{change}
``Because a single depth level gives an incomplete picture of ventilation, we also examine ideal age throughout the global ocean (Fig. 7). Averaged below 2000 m, the ageing is strongly basin dependent. The Atlantic increases only modestly, from 377 years in PI to 439 years in LGM and 508 years in MIS3, because it remains ventilated from the north. The Southern Ocean more than doubles, from 381 to 789 and 844 years, and the Pacific, already the oldest basin in the interglacials at 847 years, reaches 1291 and 1314 years. The zonal-mean sections (Fig. 7f-j) show that the glacial ageing is concentrated below roughly 2000 m and is largest in the deep Pacific, while the upper kilometre changes comparatively little. The contrast between a mildly ageing Atlantic and a strongly ageing Southern Ocean and Pacific is consistent with the abyssal cell weakening while the upper cell is maintained.''
\end{change}

\comment{R3}{Minor comments}

\comment{R3}{MH, LIG are not abbreviation that commonly used so I would remove them from the text and spell out each time. Also I don't think these were not defined at any point of the manuscript.}

\response The abbreviations are now defined at first use in the Results as well as in the Methods. We have retained them thereafter because the figures are organised in five columns labelled by these codes. We spell the names \hl{(i.e., MH, LIG) out} .

\comment{R3}{L19: 'in the model' needs to be clearly stated as this is not how AABW forms in the real modern > ocean.}

\response Added to the abstract, which now reads ``In the model, interglacial dense water formation is driven mainly by turbulent heat loss...''. The same qualification has been added in the Results.

\comment{R3}{L25: change deep from “clear imprint on deep-ocean ventilation” to abyssal Southern Ocean. As it is only one depth plot at 4000m has been shown. The deep ocean is usually considered as all the water below 1000m. Plots of overturning/ideal age globally would be helpful to assess the overall impact on ventilation.}

\response Thanks for the reviewer’s suggestion. \hl{Abyssal overturning cell and abyssal water specifically,} and the new global age figure covers the whole water column rather than a single level.

\comment{R3}{L102-105: a slightly expanded intro on the SAM may be helpful.}

\response We expanded the SAM in the introduction.

\comment{R3}{L 125: typo, more.}

\response Corrected.

\comment{R3}{L129: it would be much more useful to show the salinity minus the whole ocean change applied, to > see what arises from local dynamics.}

\response This is a good suggestion. The glacial salinity anomalies in Fig. 1 include the uniform global increase applied at initialisation to represent water stored in the ice sheets, which obscures the locally generated signal. We have noted this explicitly in the text so that readers can interpret the panels correctly.

\comment{R3}{L134: 400m is not very deep given you are comparing to what is happening to ideal age at 4000 m.}

\response We agree. The 400 m contour marks active convection rather than the depth reached by the resulting water. We now also report the 600 m statistics and the maximum mixed-layer depth, and the new overturning figure connects the surface signal to the abyssal circulation directly rather than by implication.

\comment{R3}{L144: the westerlies appear to increase under both glacial and last interglacial forcings, why?}

\response We answered this the above discussion.

\comment{R3}{L152: it needs to be clearer that this is in the model.}

\response Rephrased to ``the model produces deep mixed layers...''.

\comment{R3}{L175: I wonder if this result isn't more applicable to AAIW/SAMW formation?}

\response The reviewer's intuition is borne out by the density-class analysis described above. Our thermally dominated transformation maximum sits at γ\textsubscript{n} ≈ 27.3 kg m\textsuperscript{-3}, which is intermediate and mode water density rather than bottom water density. We now state in the Results that the interglacial thermal pathway is more relevant to intermediate and mode water formation than to bottom water proper. We are grateful for the suggestion, which turned out to be one of the more useful reframings in the revision.

\comment{R3}{L195: in the real ocean freshwater input from icesheet melting is important today.}

\response Added, with the explicit acknowledgement that this configuration has no ice shelf cavities and applies no prescribed meltwater flux, so that pathway is absent from our simulations.

In the revised manuscript this now reads:

\begin{change}
``We note that this statement applies to the freshwater sources represented in our configuration. In the present-day ocean, meltwater derived from the Antarctic ice sheet, both as basal melt beneath ice shelves and as iceberg discharge, is an important part of the coastal freshwater budget and acts to suppress dense water formation . Because the model has no ice shelf cavities and no prescribed meltwater flux, that term is absent here, and the freshwater budget we decompose is dominated by sea ice thermodynamics.''
\end{change}

\comment{R3}{L206-209: should be moved to the introduction maybe?}

\response We moved it.

\comment{R3}{L291-297: this point seems fundamental to the rest of the manuscript.}

\response Agreed. We have addressed at length in the above.

\comment{R3}{L302-305: can't you just compare you model results directly in the regions with observations?}

\response Done. See the above response and Fig. 8.

\comment{R3}{L339: this would be great, but sed rates around the Antarctic margin are very low.}

\response We agree with the reviewer. Reviewer 1 asked for a supporting reference for the same passage; \hl{we concluded the proposal was not well founded and have withdrawn it.}

\comment{R3}{L347: not the water, the carbon in the water – given the substantial and variable preformed 14C > aging.}

\response Thanks for pointing out this. We corrected it. The text now refers to the radiocarbon age of dissolved inorganic carbon.

\comment{R3}{General comment on the figures: would it be possible to show whats happening under the ice > shelves? At the moment they are just whited out. Anyway to show this?}

\response Unfortunately the model configuration does not include ice shelf cavities, so there is no ocean beneath the floating ice shelves to show. The white areas are genuinely outside the model domain rather than masked output. We have stated this in the Methods and figure captions so that readers are not left wondering, and we list it among the limitations, since the absence of cavities also means basal melt is not represented.

\comment{R3}{Please be consistent across figures and label MH, LIG, LGM, MIS3 on top of all panels.}

\response We have done throughout the manuscript. Reviewer 1 also made the same request.

%=====================================================================

% docx heading: Response to comments by Reviewer #4

\reviewer{Reviewer 4}

\comment{R4}{I co-reviewed this manuscript with one of the reviewers who provided the listed reports.}

\response We thank Reviewer 4 for the contribution to the review, and see corresponding responses.

%=====================================================================

% docx heading: Editorial requirements

\reviewer{Editorial requirements}

\response \textbf{Colour vision deficiency.} The new figures use the Okabe-Ito colour-blind-safe palette, and the transformation curves are additionally distinguished by line style. We have replaced the red/green pairing in the transformation figures with vermillion and bluish-green, and the overturning anomaly panels use a purple-orange diverging scale rather than red-green.

\textbf{Data availability.} A statement has been added. The model output needed to reproduce every figure has been assembled as a repository deposit, comprising the climatological monthly means for the five experiments, the 100-year monthly transformation time series, the Southern Annular Mode composite fields, and the ocean mesh files. The DOI will be inserted on acceptance. Source data for the line-graph figures are provided with the paper.

\textbf{Code availability.} A statement has been added. The analysis and plotting scripts, covering the surface buoyancy flux decomposition, the transformation calculations including the new sea ice sector comparison, the compositing, and all figures, are archived with the data deposit and are also available at https://github.com/xshi-awi/aabw\textsubscript{5}exps.

\textbf{Source data.} Provided for the line-graph figures, as a single spreadsheet with one sheet per figure panel.

\textbf{ORCID.} The corresponding author has linked an ORCID to the submission system, and co-authors have been asked to do the same.

\end{document}
